\begin{examplebox}
\textbf{\textcolor{CExample}{例 1｜基础：求交点与弦中点}}
\textbf{题目：}抛物线 $y^2=4x$ 上的 $A(1,2)$、$B(4,4)$，过两点的切线交于 $M$，弦中点为 $N$，求 $M,N$ 并判断 $MN$ 的位置。

\textbf{解：}由 $2p=4$ 得 $p=2$。先算中点：
\[
N\left(\frac{1+4}{2},\frac{2+4}{2}\right)
 =\left(\frac52,3\right).
\]
已知两切点纵坐标为 $2,4$，所以 $y_M=(2+4)/2=3$。
代入 $A$ 处切线 $2y=2(x+1)$：
\[
2\times3=2(x_M+1)
\ \Longrightarrow\ x_M=2.
\]
于是 $M(2,3)$，$N(5/2,3)$，两点纵坐标相同且不同点。
因为共同纵坐标为 $3\ne0$，所以 $MN$ 与 $x$ 轴\textbf{平行}，且 $|MN|=1/2$。

\medskip
\par\noindent\begin{minipage}{\linewidth}
\begin{center}
\begin{tikzpicture}[x=1.05cm,y=1.05cm,>=stealth,font=\small]
\path[use as bounding box] (-2.5,-2.7) rectangle (7.5,3.8);
\begin{scope}
\node at (0,3.3) {对称切点：$MN$ 与 $x$ 轴重合};
\draw[->,gray] (-2,0)--(2.2,0) node[right] {$x$};
\draw[->,gray] (0,-2.4)--(0,2.5) node[above] {$y$};
\draw[blue,thick,domain=-2.3:2.3,samples=70] plot ({\x*\x/4},\x);
\draw[orange,thick] (-1.4,-.4)--(1.3,2.3) (-1.4,.4)--(1.3,-2.3);
\draw[gray,thick] (1,-2)--(1,2);
\draw[red,very thick] (-1,0)--(1,0);
\fill (1,2) circle (2pt) node[right] {$A$};
\fill (1,-2) circle (2pt) node[right] {$B$};
\fill[red] (-1,0) circle (2pt) node[below left] {$M$};
\fill[red] (1,0) circle (2pt) node[above right] {$N$};
\end{scope}
\begin{scope}[xshift=5.7cm]
\node at (0,3.3) {顶点切线：$MN\parallel x$ 轴};
\draw[->,gray] (-1.5,0)--(1.8,0) node[right] {$x$};
\draw[->,gray] (0,-2.4)--(0,2.5) node[above] {$y$};
\draw[blue,thick,domain=-2.3:2.3,samples=70] plot ({\x*\x/4},\x);
\draw[orange,thick] (0,-2.2)--(0,2.3) (-1.3,-.3)--(1.3,2.3);
\draw[gray,thick] (0,0)--(1,2);
\draw[red,very thick] (0,1)--(.5,1);
\fill (0,0) circle (2pt) node[below left] {$A=O$};
\fill (1,2) circle (2pt) node[right] {$B$};
\fill[red] (0,1) circle (2pt) node[left] {$M$};
\fill[red] (.5,1) circle (2pt) node[right] {$N$};
\node[orange,left] at (0,-1.5) {$x=0$};
\end{scope}
\end{tikzpicture}


左图对应例 2，右图对应例 3；均取 $y^2=4x$，同尺度比较。
\end{center}
\end{minipage}\par

\textbf{\textcolor{CExample}{例 2｜特殊情形：两切点关于 $x$ 轴对称}}
\textbf{题目：}抛物线 $y^2=4x$ 上取 $A(1,2)$、$B(1,-2)$，求 $M,N$。

\textbf{解：}仍有 $p=2$。根据完整坐标公式，
\[
M\left(\frac{2\times(-2)}{4},\frac{2+(-2)}2\right)=(-1,0),
\qquad N\left(1,0\right).
\]
故 $|MN|=2$，而 $M,N$ 均在 $x$ 轴上，所以直线 $MN$ \textbf{就是 $x$ 轴}。
\textbf{教学提醒：}不能把“与坐标轴重合”不加区别地写成“两条不同直线平行”。

\medskip
\textbf{\textcolor{CExample}{例 3｜边界情形：一个切点恰是顶点}}
\textbf{题目：}抛物线 $y^2=4x$ 上取 $A(0,0)$、$B(1,2)$，求切线交点 $M$ 与中点 $N$。

\textbf{解：}顶点 $A$ 处切线由 $0\cdot y=2(x+0)$ 得 $x=0$，是一条竖直直线，\textbf{不需要计算斜率}。
$B$ 处切线由 $2y=2(x+1)$ 得 $y=x+1$。
两式联立得 $M(0,1)$，而
\[
N\left(\frac{0+1}{2},\frac{0+2}{2}\right)
 =\left(\frac12,1\right).
\]
故 $MN$ 水平，$|MN|=1/2$；这说明顶点情形也符合结论。

\medskip
\textbf{\textcolor{CExample}{例 4｜对称形式：$x^2=2py$}}
\textbf{题目：}抛物线 $x^2=4y$ 上 $A(2,1)$、$B(-2,1)$，求 $M,N$ 并判断 $MN$。

\textbf{解：}$p=2$，过两点的切线分别为
\[
2x=2(y+1),\qquad -2x=2(y+1).
\]
即 $x=y+1$ 与 $-x=y+1$。联立得到 $x=0,y=-1$，故 $M(0,-1)$。
中点
\[
N\left(\frac{2+(-2)}2,\frac{1+1}2\right)=(0,1).
\]
所以 $MN$ 是\textbf{$y$ 轴本身}，其长度为 $2$；这也是“重合”的典型情形。

\medskip
\textbf{\textcolor{CExample}{例 5｜综合：斜率与 $MN$ 长度确定弦}}
\textbf{题目：}抛物线 $y^2=2x$ 上有不同两点 $A,B$，切线交于 $M$，弦中点为 $N$。
已知 $AB$ 的斜率为 $1$，$|MN|=2$，求直线 $AB$ 的方程。

\textbf{解：}
\begin{description}[style=nextline,leftmargin=2.8em,labelsep=0.5em,itemsep=0.4em]
\item[\textbf{第 1 步：由斜率求纵坐标之和。}]
这里 $p=1$，$x_i=y_i^2/2$。
由于斜率已知且等于 $1$，$x_2-x_1\ne0$，并且
\[
\begin{aligned}
1=k_{AB}
&=\frac{y_2-y_1}{x_2-x_1}\\
&=\frac{y_2-y_1}{(y_2^2-y_1^2)/2}\\
&=\frac{2}{y_1+y_2}.
\end{aligned}
\]
故 $y_1+y_2=2$，中点的纵坐标为 $y_N=1$。

\item[\textbf{第 2 步：由 $MN$ 长度求纵坐标之差。}]
本题 $p=1$，所以
\[
2=|MN|=\frac{(y_1-y_2)^2}{4}
\quad\Longrightarrow\quad (y_1-y_2)^2=8.
\]
把“和的平方”与“差的平方”相加：
\[
2(y_1^2+y_2^2)=(y_1+y_2)^2+(y_1-y_2)^2=4+8=12.
\]
因此 $y_1^2+y_2^2=6$，
\[
x_N=\frac{x_1+x_2}2=\frac{y_1^2+y_2^2}{4}=\frac32.
\]
故 $N(3/2,1)$。

\item[\textbf{第 3 步：用点斜式确定弦方程。}]
已知 $AB$ 过 $N(3/2,1)$ 且斜率为 $1$，
\[
y-1=x-\frac32
\quad\Longrightarrow\quad\boxed{y=x-\frac12}.
\]
\end{description}
\textbf{检查：}将 $x=y+1/2$ 代入 $y^2=2x$ 得 $y^2-2y-1=0$，
两根为 $1\pm\sqrt2$，确实不同；故所求直线与抛物线有两个不同交点。
\end{examplebox}
